\section{Quy trình dữ liệu}

\subsection{Thu thập dữ liệu}

\subsubsection{Lựa chọn nguồn dữ liệu}

Nhóm đã lựa chọn hai sàn thương mại điện tử \textbf{Tiki} và \textbf{eBay} làm nguồn dữ liệu chính vì các lý do sau:
\begin{itemize}
    \item \textbf{Đại diện đa dạng cho thị trường (Nội địa \& Quốc tế):} Tiki là sàn thương mại điện tử uy tín tại Việt Nam, trong khi eBay là nền tảng toàn cầu. Sự kết hợp này cho phép đối chiếu hành vi tiêu dùng, chính sách vận chuyển và chiến lược định giá ở hai quy mô thị trường khác nhau.
    \item \textbf{Độ phủ dữ liệu rộng:} Thu thập hàng các sản phẩm trải dài trên nhiều danh mục trọng điểm (Thiết bị số, Thời trang, Mỹ phẩm, Đồ gia dụng,...).
\end{itemize}

\subsubsection{Phương pháp thu thập}

\textbf{Kỹ thuật sử dụng:} Thu thập qua RESTful API kết hợp cơ chế quản lý tiến trình (Checkpoint \& Resume) và phân trang tự động.

\textbf{Công cụ và thư viện chính:}
\begin{itemize}
    \item \texttt{requests} \& \texttt{urllib3.util.retry}: Gửi HTTP requests kết hợp \texttt{HTTPAdapter} để tự động thử lại (retry) khi gặp lỗi mạng hoặc máy chủ quá tải (HTTP 429, 500, 502), đảm bảo pipeline không bị gián đoạn.
    \item \texttt{base64} \& \texttt{json}: Xử lý mã hóa OAuth 2.0 (eBay) và lưu trữ trạng thái cào dữ liệu (Tiki Checkpoint).
\end{itemize}

\textbf{Quy trình thu thập (Pipeline):}
\begin{enumerate}
    \item \textbf{Khởi tạo tham số \& xác thực:} Khai báo danh sách từ khóa/danh mục mục tiêu. Khởi tạo phiên làm việc với các Headers chống giả mạo User-Agent (Tiki) hoặc gọi API đổi Access Token qua cơ chế Client Credentials (eBay).
    \item \textbf{Quét danh sách \& phân bổ:} Phân bổ số lượng sản phẩm mục tiêu đều cho từng từ khóa. Dùng vòng lặp phân trang (\texttt{page} hoặc \texttt{offset}) để gom hàng loạt ID sản phẩm.
    \item \textbf{Truy xuất chi tiết:} Đối với Tiki, sử dụng các ID đã quét để gọi tiếp vào endpoint Detail (\texttt{/api/v2/products/\{id\}}) nhằm lấy dữ liệu đặc tả toàn vẹn nhất.
\end{enumerate}

\subsection{Thiết kế cấu trúc dữ liệu}

\subsubsection{Mô hình dữ liệu thô}

Dữ liệu thô được thu thập trực tiếp từ API và lưu vào hai file CSV. Mỗi trường phản ánh đúng response của API nguồn — chưa qua bất kỳ biến đổi nào.

\textit{(a) Tiki} — \texttt{data/raw/products.csv} (11 trường):
\begin{itemize} [leftmargin=1.5cm]
    \item \texttt{id} [INT]: định danh sản phẩm duy nhất trên Tiki.
    \item \texttt{name} [STR]: tên đầy đủ của sản phẩm.
    \item \texttt{price} / \texttt{original\_price} [INT, VND]: giá bán hiện tại và giá gốc trước chiết khấu.
    \item \texttt{discount\_rate} [INT, \%]: tỷ lệ chiết khấu do Tiki tính toán.
    \item \texttt{rating\_average} [FLOAT, 0--5] / \texttt{review\_count} [INT]: điểm đánh giá trung bình và số lượt review.
    \item \texttt{quantity\_sold} [INT]: tổng số đơn hàng đã bán (giá trị tích lũy).
    \item \texttt{brand} [STR]: nhãn hàng; có thể rỗng nếu người bán không khai báo.
    \item \texttt{category} [STR]: danh mục cấp 1 của sản phẩm.
    \item \texttt{url\_key} [STR]: slug URL dùng để tái truy xuất trang sản phẩm.
\end{itemize}

\textit{(b) eBay} — \texttt{data/raw/ebay\_products.csv} (24 trường chính):
\begin{itemize} [leftmargin=1.5cm]
    \item \texttt{item\_id} / \texttt{legacy\_item\_id} [STR]: định danh listing trên eBay.
    \item \texttt{title} [STR]: tiêu đề listing do người bán đặt.
    \item \texttt{price} [FLOAT, USD] / \texttt{currency} [STR]: giá niêm yết và đơn vị tiền tệ.
    \item \texttt{shipping\_cost} [FLOAT, USD] / \texttt{has\_free\_shipping} [BOOL]: phí và chính sách vận chuyển.
    \item \texttt{total\_cost} [FLOAT, USD]: tổng chi phí (price + shipping\_cost).
    \item \texttt{condition} [STR]: tình trạng sản phẩm (New, Used, Refurbished,...).
    \item \texttt{buying\_options} [STR]: hình thức mua (Fixed Price, Auction,...).
    \item \texttt{seller\_username} [STR] / \texttt{seller\_feedback\_score} [INT] / \texttt{seller\_feedback\_percent} [FLOAT]: định danh và uy tín người bán.
    \item \texttt{category\_id} [STR] / \texttt{category\_path} [STR]: danh mục sản phẩm.
    \item \texttt{item\_location\_country} [STR]: quốc gia của người bán.
    \item \texttt{item\_creation\_date} [STR]: ngày đăng listing.
    \item \texttt{query} [STR]: từ khóa tìm kiếm dùng để thu thập bản ghi này.
\end{itemize}

\subsubsection{Quy mô dữ liệu thô}
Ta có quy mô dữ liệu như sau:
\begin{itemize} [leftmargin=1.5cm]
    \item \texttt{data/raw/products.csv}: $\sim$10,000 sản phẩm, 11 cột, 2.3 MB — thu thập từ Tiki tháng 4/2026.
    \item \texttt{data/raw/ebay\_products.csv}: $\sim$12,000 sản phẩm, 24 cột, 9.5 MB — thu thập từ eBay tháng 4/2026.
\end{itemize}

\textbf{Tổng: $\sim$22,000 bản ghi thô, đa danh mục — đáp ứng yêu cầu hơn 5,000 dòng theo đề bài.}


\subsection{Tiền xử lý dữ liệu}

\subsubsection{Kiểm tra tính nhất quán}

\textbf{Chiến lược xử lý:}
\begin{itemize}
    \item Toàn bộ tập dữ liệu: Kiểm tra dòng trùng lặp và dòng rỗng hoàn toàn. Dữ liệu nguyên vẹn nên không cần xóa dòng.
    \item Cột \texttt{subtitle}, \texttt{item\_end\_date} (eBay): Xóa cột do tỷ lệ thiếu hụt dữ liệu quá lớn (trên 90\%), chỉ gây nhiễu không gian đặc trưng.
    \item Các cột phân loại/danh xưng (Categorical - eBay): Điền giá trị xuất hiện nhiều nhất để bảo toàn nguyên vẹn tính danh xưng và tránh tạo ra các mối quan hệ thứ tự giả tạo.
    \item Các cột định lượng hoặc biến đếm (Tiki/eBay): Điền giá trị trung vị hoặc số 0 để đề phòng lỗi khuyết dữ liệu mà không bị ảnh hưởng bởi các giá trị cực đoan.
    \item Cột \texttt{brand}, \texttt{category} (Tiki) và \texttt{category\_path} (eBay): Chuẩn hóa chuỗi bằng cách cắt bỏ khoảng trắng thừa và thay thế các chuỗi rỗng giả mạo (`Nan', `None', `Root', `Oem') thành `Unknown'.
\end{itemize}

\subsubsection{Phát hiện và xử lý ngoại lai}

\textbf{Phương pháp phát hiện:}
\begin{itemize}
    \item Sử dụng IQR để xác định bằng con số cụ thể ranh giới ngoại lai, mang lại tính bền vững cao hơn Z-score đối với dữ liệu E-commerce có phân phối lệch cực mạnh.
    \item Trực quan hóa bằng biểu đồ Boxplot kết hợp với Histogram sử dụng kỹ thuật "Zoom-in" để quan sát rõ cấu trúc phân phối cốt lõi.
\end{itemize}

\textbf{Chiến lược xử lý:} 
 Thay vì xóa bỏ dữ liệu, hệ thống chỉ loại bỏ 1\% các giá trị tài chính cao vô lý tại các cột \texttt{price}, \texttt{original\_price} (Tiki) và \texttt{price}, \texttt{shipping\_cost} (eBay). Các dòng dữ liệu có chỉ số ngoại lai thuộc về tương tác lõi như \texttt{quantity\_sold} hay  \texttt{seller\_feedback\_score} được giữ nguyên tuyệt đối nhằm bảo toàn nhóm sản phẩm Best-seller và các nhà bán hàng uy tín nhất.

\subsubsection{Chuyển đổi kiểu dữ liệu}

\textbf{Chiến lược chuyển đổi:}
\begin{itemize}
    \item Cột \texttt{id} (Tiki) và \texttt{item\_id} (eBay): Ép kiểu sang chuỗi (\texttt{String/Object}) để bảo vệ các trường định danh, tránh việc thuật toán tự động thực hiện các phép toán thống kê sai lệch.
    \item Cột \texttt{price}, \texttt{original\_price} và \texttt{shipping\_cost}: Ép kiểu đồng nhất sang định dạng số (\texttt{int64/float64}) để sẵn sàng cho các phép toán thống kê.
    \item Cột \texttt{item\_creation\_date}, \texttt{item\_end\_date} (eBay): Ép kiểu sang định dạng thời gian chuẩn \texttt{datetime64[ns, UTC]} để phục vụ tính toán chuỗi thời gian và vòng đời sản phẩm.
\end{itemize}

\subsection{Kết quả sau tiền xử lý}

\begin{table}[H]
    \centering
    \caption{Thống kê dữ liệu sau tiền xử lý}
    \begin{tabular}{|l|c|c|}
        \hline
        \textbf{Chỉ số} & \textbf{Trước xử lý} & \textbf{Sau xử lý} \\
        \hline
        Tổng số dòng & [X] & [Y] \\
        \hline
        Missing values & [X]\% & 0\% \\
        \hline
        Outliers & [X] dòng & [Y] dòng \\
        \hline
        Duplicate rows & [X] dòng & 0 dòng \\
        \hline
    \end{tabular}
\end{table}

\clearpage
